\documentclass[10pt,a4paper]{article}

% ---------- Packages ----------
\usepackage[a4paper,top=0.44in,bottom=0.44in,left=0.48in,right=0.48in]{geometry}
\usepackage{enumitem}
\usepackage{hyperref}
\usepackage{titlesec}
\usepackage{xcolor}

% ---------- Colors & Links ----------
\definecolor{linkblue}{HTML}{0D47A1}
\hypersetup{
    colorlinks=true,
    linkcolor=black,
    urlcolor=linkblue
}

% ---------- Formatting & Spacing ----------
\setlist[itemize]{noitemsep, topsep=1.5pt, itemsep=1.2pt, parsep=0pt, partopsep=0pt, leftmargin=1.15em}
\setlength{\parindent}{0pt}
\setlength{\parskip}{0pt}

\titleformat{\section}{\large\bfseries}{}{0em}{}[\vspace{-2pt}\titlerule]
\titlespacing*{\section}{0pt}{5pt plus 1pt minus 1pt}{2.5pt plus 1pt minus 1pt}

\begin{document}
\pagestyle{empty}

% ---------- Header ----------
\begin{center}
    {\LARGE \textbf{Shahriar Alom Masud}} \\[3pt]
    \small
    Dhaka, Bangladesh \enspace$|$\enspace +880 1740-071118 \enspace$|$\enspace \href{mailto:shahriar0002@std.uftb.ac.bd}{shahriar0002@std.uftb.ac.bd} \\[1pt]
    \href{https://www.linkedin.com/in/shahriar-alom-masud}{LinkedIn} \enspace$|$\enspace \href{https://github.com/Masud744}{GitHub} \enspace$|$\enspace \href{https://masud-hj9g.onrender.com/}{Portfolio}
\end{center}
\vspace{-3pt}

% ---------- Professional Summary ----------
\section*{Professional Summary}
Embedded Software and IoT Engineer with 2+ years of hands-on experience developing C/C++ firmware for microcontrollers. Skilled in FreeRTOS, dual-partition OTA updates, wireless sensor communication, and running lightweight machine learning models on edge devices.

% ---------- Skills ----------
\section*{Skills}
\textbf{Languages:} C, C++, Python, JavaScript, TypeScript, Dart, SQL, Bash \\[1pt]
\textbf{Embedded:} ESP32, ESP8266, Raspberry Pi, Arduino, FreeRTOS, Dual-Partition OTA, NVS, EEPROM, Hardware Timers, ISR \\[1pt]
\textbf{Protocols:} MQTT, HTTP/HTTPS, LoRa, BLE, UART, SPI, I2C, 1-Wire, WebSockets \\[1pt]
\textbf{AI / ML:} TinyML, ONNX Runtime, OpenCV, YOLOv8, PyTorch, Scikit-learn, MediaPipe \\[1pt]
\textbf{Cloud \& Tools:} Azure, FastAPI, Supabase, Firebase, MongoDB, InfluxDB, Docker, React, Flutter, Git, GitHub Actions, PlatformIO, Linux

% ---------- Experience ----------
\section*{Experience}

\textbf{Open Source Developer --- SmartProv} \hfill \textit{Apr 2026 -- Present} \\
\href{https://github.com/Masud744/SmartProv}{github.com/Masud744/SmartProv} \,|\, \textit{Arduino Library Manager, ESP32, ESP8266, C++, NVS/EEPROM}
\begin{itemize}
    \item Built and published \textbf{SmartProv} on the Arduino Library Manager --- an open-source Wi-Fi provisioning library for ESP32 and ESP8266 that lets users configure device credentials through a captive portal instead of reflashing firmware.
    \item Stores Wi-Fi credentials in NVS and EEPROM with multi-network fallback and factory-reset support, so devices reconnect automatically after power loss without manual reflashing.
    \item Cuts device network setup time from several minutes to under 30 seconds per node, removing the need to open Arduino IDE for every Wi-Fi change.
\end{itemize}

% ---------- Projects ----------
\section*{Projects}

\textbf{RespiGuard.ai --- Cyber-Physical Edge AI Respiratory Monitoring} \hfill \textit{May 2026 -- Sep 2026} \\
\href{https://github.com/Masud744/RespiGuard}{github.com/Masud744/RespiGuard} \,|\, \href{https://respiguard-fvh9.onrender.com/}{Live: respiguard-fvh9.onrender.com} \\
\textit{ESP32, FreeRTOS, C/C++, PMS5003 (UART), MQ-135 (ADC), DHT22, SSD1306 OLED, FastAPI, TreeSHAP, React}
\begin{itemize}
    \item Built a multi-sensor ESP32 edge node reading laser particulate matter ($PM_{1.0}/2.5/10$ via UART2), calibrated toxic gas (MQ-135 via 12-bit ADC1), and climate data using FreeRTOS tasks and SSD1306 I2C display.
    \item Implemented dual-pipeline edge/clinical inference with TreeSHAP explainability to evaluate asthma risk zones (Green, Yellow, Red based on PEFR), streaming HMAC-SHA256 authenticated telemetry over HTTPS to a FastAPI and React dashboard.
\end{itemize}
\vspace{2.5pt}

\textbf{Smart Solar-Grid HEMS --- Intelligent Load Scheduling \& Power Routing} \hfill \textit{Jan 2026 -- Sep 2026} \\
\href{https://github.com/Masud744/solar-aware-hems}{github.com/Masud744/solar-aware-hems} \,|\, \href{https://solar-aware-hems.vercel.app}{Live: solar-aware-hems.vercel.app} \\
\textit{ESP32, C++, FreeRTOS, ZMPT101B, ACS712, 8-Relay Matrix, NVS Flash, Python, FastAPI, Scikit-learn, React}
\begin{itemize}
    \item Engineered dual-core ESP32 firmware running a 1.5--2.0\,kHz True-RMS AC sampling loop on Core 1 (ZMPT101B voltage, ACS712 current with NVS calibration) and Wi-Fi cloud telemetry on Core 0.
    \item Designed an 8-relay dual-bank matrix (Grid vs Solar) with 300\,ms break-before-make dead-time protection and anti-chatter hysteresis, scheduling loads using a $k \times \sigma$ risk-aware solar forecasting engine.
\end{itemize}
\vspace{2.5pt}

\textbf{ESP32 GitHub OTA --- Fail-Safe Firmware Update \& Recovery System} \hfill \textit{Jun 2026 -- Aug 2026} \\
\href{https://github.com/Masud744/esp32-smartprov-github-ota}{github.com/Masud744/esp32-smartprov-github-ota} \\
\textit{ESP32, C/C++, Dual-Partition Flash (ota\_0/ota\_1), NVS, HTTPS, GitHub REST, PlatformIO}
\begin{itemize}
    \item Developed a serverless OTA client on ESP32 that checks GitHub-hosted \texttt{version.json} over HTTPS and streams compiled binaries into inactive flash partitions with automatic download retry limits.
    \item Built crash-recovery logic using NVS boot counters; automatically triggers \texttt{Update.rollBack()} if a new build crashes 3 times or drops Wi-Fi, blacklisting faulty releases to prevent infinite boot-loops.
\end{itemize}

% ---------- Education ----------
\section*{Education}

\textbf{University of Frontier Technology, Bangladesh} \hfill \textit{Aug 2023 -- Jun 2027 (Expected)} \\
Bachelor of Science (Engineering) in Internet of Things \& Robotics Engineering \hfill \textit{Gazipur, Bangladesh}

% ---------- Leadership ----------
\section*{Leadership}

\textbf{Head of Membership Development --- UFTB Robotics Club} \hfill \textit{Oct 2025 -- Present} \\
Organized hands-on technical workshops on microcontroller interfacing, C/C++ firmware debugging, and sensor-to-cloud telemetry for 60+ junior engineering students. \\[2.5pt]
\textbf{Robo Soccer Segment Leader --- RoboFusion 1.0} \hfill \textit{Aug 2026} \\
Led and coordinated the Robo Soccer segment across 20 participating inter-university teams, overseeing arena technical setups, wireless frequency allocations, and match execution.

\end{document}
